% The Slide Definitions
%document
\documentclass[10pt]{beamer}
%theme
\usetheme{metropolis}
% packages
\usepackage{color}
\usepackage{listings}
\usepackage[ngerman]{babel}
\usepackage[utf8]{inputenc}
\usepackage{multicol}


% color definitions
\definecolor{mygreen}{rgb}{0,0.6,0}
\definecolor{mygray}{rgb}{0.5,0.5,0.5}
\definecolor{mymauve}{rgb}{0.58,0,0.82}

\lstset{
    backgroundcolor=\color{white},
    % choose the background color;
    % you must add \usepackage{color} or \usepackage{xcolor}
    basicstyle=\footnotesize\ttfamily,
    % the size of the fonts that are used for the code
    breakatwhitespace=false,
    % sets if automatic breaks should only happen at whitespace
    breaklines=true,                 % sets automatic line breaking
    captionpos=b,                    % sets the caption-position to bottom
    commentstyle=\color{mygreen},    % comment style
    % deletekeywords={...},
    % if you want to delete keywords from the given language
    extendedchars=true,
    % lets you use non-ASCII characters;
    % for 8-bits encodings only, does not work with UTF-8
    frame=single,                    % adds a frame around the code
    keepspaces=true,
    % keeps spaces in text,
    % useful for keeping indentation of code
    % (possibly needs columns=flexible)
    keywordstyle=\color{blue},       % keyword style
    % morekeywords={*,...},
    % if you want to add more keywords to the set
    numbers=left,
    % where to put the line-numbers; possible values are (none, left, right)
    numbersep=5pt,
    % how far the line-numbers are from the code
    numberstyle=\tiny\color{mygray},
    % the style that is used for the line-numbers
    rulecolor=\color{black},
    % if not set, the frame-color may be changed on line-breaks
    % within not-black text (e.g. comments (green here))
    stepnumber=1,
    % the step between two line-numbers.
    % If it's 1, each line will be numbered
    stringstyle=\color{mymauve},     % string literal style
    tabsize=4,                       % sets default tabsize to 4 spaces
    % show the filename of files included with \lstinputlisting;
    % also try caption instead of title
    language = Python,
	showspaces = false,
	showtabs = false,
	showstringspaces = false,
	escapechar = ,
}

\def\ContinueLineNumber{\lstset{firstnumber=last}}
\def\StartLineAt#1{\lstset{firstnumber=#1}}
\let\numberLineAt\StartLineAt



\newcommand{\codeline}[1]{
	\alert{\texttt{#1}}
}


% Author and Course information
% This Document contains the information about this course.

% Authors of the slides
\author{Felix Döring, Felix Wittwer, Anton Obersteiner}

% Name of the Course
\institute{Python-Kurs}

% Fancy Logo 
\titlegraphic{\hfill\includegraphics[height=1.25cm]{../templates/fsr_logo_cropped}}



% Custom Bindings
% \newcommand{\codeline}[1]{
%	\alert{\texttt{#1}}
%}


% Presentation title
\title{Glossar}
\date{\today}
\setlength{\columnsep}{.4cm}

\begin{document}

\maketitle

\begin{frame}{Gliederung}
	\begin{multicols}{2}
		\setbeamertemplate{section in toc}[sections numbered]
		\tableofcontents
	\end{multicols}
\end{frame}


\section{Basic}

\subsection{Variablen}
\begin{frame}{Variable}
	\lstinputlisting{resources/glossar/variables.py}
	\begin{itemize}
		\item Speichern Werte, bis sie neu gesetzt werden
		\item Variablen in einer Funktion sind nicht außerhalb sichtbar
	\end{itemize}
\end{frame}

\subsection{if-elif-else}
\begin{frame}{if-elif-else}
	\lstinputlisting{resources/glossar/if_elif_else.py}
	\begin{itemize}
		\item \texttt{if}-Bedingung wird zuerst geprüft
		\item Wenn die nicht gilt (\texttt{True} ist), wird das erste \texttt{elif} überprüft
		\item Wenn keine \texttt{if}- oder \texttt{elif}-Bedingung gilt, kann \texttt{else} ausgeführt werden
		\item \texttt{elif} und \texttt{else} muss es nicht geben
	\end{itemize}
\end{frame}

\subsection{\texttt{print} und \texttt{input}}
\begin{frame}{print und input}
	\lstinputlisting{resources/glossar/print_input.py}
	\begin{itemize}
		\item \texttt{print} gibt alle Argumente aus
		\item \texttt{input} gibt seinen prompt aus und wartet auf Nutzerinput
		\item Vorsicht: \texttt{input} gibt immer einen String zurück
	\end{itemize}
\end{frame}

\subsection{Funktionen}
\begin{frame}{Funktionen}
	\lstinputlisting{resources/glossar/function.py}
	\begin{itemize}
		\item mit \texttt{funktion(argumente)} aufrufen
		\item \texttt{return}-Wert kann weiter verwendet werden
		\item alles andere, was die Funktion macht, sind Seiteneffekte
	\end{itemize}
\end{frame}

\subsection{Datentypen}
\begin{frame}{Datentypen}
	\begin{description}
		\item[int] Ganzzahl: \texttt{5, -128, int(17.5) == 17}
		\item[float] Kommazahl: \texttt{12.7, -0.001}
		\item[str] Zeichenkette: \texttt{"Text mit 123", "q", str(5) == "5"}
		\item[bool] Ja-Nein-Wert: \texttt{a == 5} $\to$ \texttt{True} oder \texttt{False}
		\item[type] Datentyp: \texttt{type("q") == str}
		\item[list] Liste mit irgendwelchen Elementen: \texttt{["Houlle", 1, 12]}
		\item[dict] Zuordnung \texttt{key} $\to$ \texttt{element}: \texttt{\{1: "ja", 2:"nö"\}}
	\end{description}
	\lstinputlisting{resources/glossar/datatypes.py}
\end{frame}

\subsection{Operatoren}
\begin{frame}{Operatoren}
	\begin{description}
		\item[\tt+] Addiert Zahlen, hängt anderes hintereinander
		\item[\tt*] Multipliziert Zahlen, vervielfältigt anderes
		\item[\tt-] Subtrahiert Zahlen
		\item[\tt/] Teilt Zahlen
		\item[\tt\%] Bildet den Rest $a \mod b$
		\item[\tt**] Exponent: statt \^{}
		\item[\tt==] Sind sie gleich? $\to$ bool als Ergebnis
	\end{description}
	\lstinputlisting{resources/glossar/operators.py}
\end{frame}
\begin{frame}{Operatoren}
	\begin{description}
		\item[\tt==] gleich?
		\item[\tt!=] ungleich?
		\item[\tt<, >] kleiner? größer?
		\item[\tt<=, >=] kleiner gleich? größer gleich
		\item[\tt and] gelten $a$ und $b$?
		\item[\tt or] gelten $a$ und/oder $b$?
	\end{description}
	\lstinputlisting{resources/glossar/operators_bool.py}
\end{frame}

\subsection{Listen}
\begin{frame}{Listen}
	\begin{itemize}
		\item Mehrere Werte mit Reihenfolge
		\item Index \texttt{0, 1, 2, ..., len(liste) - 1}
	\end{itemize}
	\lstinputlisting{resources/glossar/list_squares.py}
\end{frame}

\subsection{\texttt{for}-Loop und \texttt{range}}
\begin{frame}{\texttt{for}-Loop und \texttt{range}}
	\begin{itemize}
		\item Syntax: \texttt{for $a$ in $A$: $block$}
		\item $a$ durchläuft alle Elemente von $A$ einmal
		\item jedes mal wird $block$ (Eingerücktes darunter) ausgeführt
		\item \texttt{range(start = 0, stop, step = 1)} ist wie eine Liste von \texttt{start} bis \texttt{stop}
		\item Enden von Bereichen (hier: \texttt{stop}) werden nicht mit berechnet!
		\item eins größer wählen als letztes Element
	\end{itemize}
	\lstinputlisting{resources/glossar/for_range.py}
\end{frame}

\subsection{\texttt{while}-Loop}
\begin{frame}{\texttt{while}-Loop}
	\begin{itemize}
		\item Syntax: \texttt{while $b$: $block$}
		\item prüft Bedingung $b$ immer wieder
		\item führt jedes Mal, wenn $b$ gilt, $block$ aus
	\end{itemize}
	\lstinputlisting{resources/glossar/while.py}
\end{frame}

\subsection{\texttt{continue} und \texttt{break}}
\begin{frame}{\texttt{continue} und \texttt{break}}
	\begin{description}
		\item[\texttt{continue}] Schleife ignoriert den Rest des Blocks
		\item[]	und geht zum nächsten Durchlauf über
		\item[\texttt{break}] Schleife bricht komplett ab
	\end{description}
\end{frame}

\section{Intermediate}

\subsection{Files}
\begin{frame}{Files}
	\lstinputlisting{resources/glossar/files.py}
	\lstinputlisting{resources/glossar/files_with.py}
\end{frame}

\section{Klassen}
\subsection{Klassen, Objekte und Attribute}
\begin{frame}{Klassen und Attribute}
	\begin{itemize}
		\item<1-> Klassen sind wie Schablonen
		\item<2-> Objekte sind konkrete Instanzen
	\end{itemize}
	\onslide<3->\lstinputlisting{resources/02_grundlagen/class.py}
	\onslide<3->{auf Deutsch: selbst festgelegt}
\end{frame}

\begin{frame}{Attribute}
	Attribute setzen: meist im \textbf{Initialisierer} \alert{\tt \_\_init\_\_}
	\lstinputlisting{resources/02_grundlagen/attributes.py}
	auf Deutsch: selbst festgelegt
\end{frame}

\subsection{Methoden}
\begin{frame}{Methoden}
	\begin{itemize}
		\item Funktionen, die Teil von Klassen/Objekten sind
		\item<2-> erstes Argument \texttt{self} wird übergeben
		\item<3-> \texttt{	def method(self, arg2): ...}
		\item<3-> \texttt{obj.method(arg2)}
	\end{itemize}
\end{frame}

\subsection{Super- und Subklassen}
\begin{frame}{Super- und Subklassen}
	Erben / Erweitern: mit \alert{\tt class subclass(superclass):} definieren
	\begin{itemize}
		\item neue Variablen und Methoden hinzufügen, auch alte überschreiben
		\item Superklasse mit \alert{\tt super()} aufrufen
	\end{itemize}
	\lstinputlisting{resources/02_grundlagen/inheritance.py}
\end{frame}
\subsection{Spezielle Methoden}
\begin{frame}{Spezielle Methoden}
	von Python intern verwendet

	\begin{description}
		\item<1->[\tt\_\_init\_\_] Bei Instanziierung aufgerufen
		\item<2->[\tt\_\_del\_\_] Bei Löschung aufgerufen (selten)
		\item<3->[\tt\_\_str\_\_] Für Ausgabe als String
		\item<4->[\tt\_\_add\_\_] Funktion für \alert{\tt +}-Operator
		\item<5->[\tt\_\_iadd\_\_] Funktion für \alert{\tt +=}-Operator (muss self zurückgeben)
		\item<6->[\tt\_\_sub\_\_] Funktion für \alert{\tt -}-Operator
		\item<7->[\tt\_\_truediv\_\_] Funktion für \alert{\tt /}-Operator
	\end{description}
\end{frame}


\end{document}
